% Introduction.tex

\chapter{مقدمه}
\label{chap:intro}

در این فصل، ابتدا در \autoref{sec:background} زمینه کلی پژوهش و اهمیت یادگیری فدرال در محیط‌های ابری چندگانه بررسی می‌شود. سپس در \autoref{sec:problem}، مسئله اصلی این پژوهش تشریح می‌گردد. در \autoref{sec:motivation}، چالش‌ها و انگیزه‌های انجام این کار بیان می‌شود و در \autoref{sec:contributions}، نوآوری‌های این پژوهش به‌اختصار فهرست می‌شود. در نهایت، \autoref{sec:structure} ساختار کلی این گزارش را توضیح می‌دهد.

\section{پیش‌زمینه}
\label{sec:background}

در دهه اخیر، آموزش مدل‌های یادگیری ماشین به یکی از محورهای اصلی توسعه فناوری اطلاعات تبدیل شده است. رویکرد سنتی در این حوزه بر پایه تمرکز داده‌ها در یک سرور مرکزی استوار بود؛ تمام داده‌های تولیدشده توسط کاربران یا سیستم‌های مختلف به یک مخزن مشترک منتقل می‌شدند و آموزش مدل بر روی همین مخزن انجام می‌گرفت. این رویکرد اگرچه از دیدگاه سادگی پیاده‌سازی جذاب به نظر می‌رسد، اما در عمل با چالش‌های جدی روبه‌رو است؛ انتقال حجم زیادی از داده در بستر شبکه هزینه‌بر است، حریم خصوصی کاربران را به خطر می‌اندازد و در بسیاری از حوزه‌ها نظیر پزشکی، مالی و نظامی با محدودیت‌های قانونی مواجه است.

یادگیری فدرال
\gls{FL}
که نخستین بار در سال ۲۰۱۷ توسط گوگل معرفی شد، رویکردی بدیل برای آموزش توزیع‌شده ارائه می‌دهد. در این رویکرد، هر شرکت‌کننده (که از این پس «گره» نامیده می‌شود) مدل را به‌صورت محلی بر روی داده‌های خود آموزش می‌دهد و تنها به‌روزرسانی‌های مدل، نه داده‌های خام، را برای تجمیع به سرور مرکزی ارسال می‌کند. این ویژگی موجب می‌شود که داده‌های حساس هرگز از محیط محلی خارج نشوند.

\begin{figure}
\centering
\includegraphics[width=0.85\linewidth]{/FedLearning/overview}
\caption{\lofimage{/FedLearning/overview}%
نمای کلی از معماری یادگیری فدرال در محیط ابر چندگانه. هر گره داده محلی خود را نگه می‌دارد و تنها به‌روزرسانی مدل را به سرور ارسال می‌کند.}
\label{fig:fl_overview}
\end{figure}

با گسترش خدمات ابری و پیدایش محیط‌های ابری چندگانه، یادگیری فدرال کاربردهای جدیدی یافته است. در چنین محیط‌هایی، سازمان‌ها و شرکت‌های مختلف از ارائه‌دهندگان ابری گوناگون نظیر
\gls{AWS}،
\gls{GCP}
و
\lr{Azure}
بهره می‌برند. هر یک از این محیط‌ها ممکن است داده‌هایی با توزیع آماری متفاوت داشته باشند و از زیرساخت شبکه‌ای با مشخصات متمایز بهره‌مند باشند. آموزش یک مدل مشترک در چنین بستری، مسائل پیچیده‌ای را پیش می‌آورد که حل آن‌ها هدف اصلی این پژوهش است.


\section{طرح مسئله}
\label{sec:problem}

در یک محیط ابری چندگانه با $N$ گره شرکت‌کننده، فرض کنید هر گره $i$ مجموعه‌ای از داده‌های محلی $\mathcal{D}_i$ را در اختیار دارد. هدف کلی این است که بدون تبادل داده‌های خام، یک مدل سراسری $w^*$ بیابیم که تابع هدف زیر را کمینه کند:

\begin{equation}
w^* = \arg\min_{w} \sum_{i=1}^{N} \frac{|\mathcal{D}_i|}{|\mathcal{D}|} \mathcal{L}_i(w; \mathcal{D}_i)
\label{eq:global_objective}
\end{equation}

که در آن $\mathcal{L}_i$ تابع خطای محلی گره $i$ و $|\mathcal{D}| = \sum_{i=1}^{N} |\mathcal{D}_i|$ اندازه کل داده‌ها است.

رویکرد استاندارد برای حل این مسئله، الگوریتم \gls{FedAvg} است که در هر دور ارتباطی، مراحل زیر را طی می‌کند: سرور مدل سراسری را به گره‌های منتخب ارسال می‌کند، هر گره چند دوره آموزش محلی انجام می‌دهد، سپس به‌روزرسانی‌های مدل به سرور بازمی‌گردند و سرور این به‌روزرسانی‌ها را با میانگین وزن‌دار تجمیع می‌کند.

اما در محیط‌های ابری واقعی، این فرض ساده با واقعیت فاصله دارد. چند دسته مشکل اساسی وجود دارد که \gls{FedAvg} استاندارد را در این بستر ناکارآمد می‌کند:

\begin{itemize}
    \tick \textbf{ناهمگونی محاسباتی}: گره‌های مختلف روی انواع متفاوتی از ماشین‌های مجازی یا نمونه‌های ابری اجرا می‌شوند. یک نمونه با \gls{GPU} قوی ممکن است ده‌ها برابر سریع‌تر از یک نمونه \lr{CPU}-محور باشد. این تفاوت موجب می‌شود که برخی گره‌ها به‌موقع پاسخ ندهند و دور ارتباطی تأخیر داشته باشد.
    
    \tick \textbf{ناهمگونی شبکه}: ارتباط بین نواحی مختلف یک ارائه‌دهنده ابری یا میان ارائه‌دهندگان مختلف، از نظر تأخیر و پهنای باند یکسان نیست. پیوند درون‌ناحیه‌ای ممکن است تأخیر کمتر از یک میلی‌ثانیه داشته باشد، در حالی که پیوند بین دو قاره می‌تواند تأخیر بیش از ۱۵۰ میلی‌ثانیه ایجاد کند.
    
    \tick \textbf{ناهمگونی داده}: هر ناحیه ابری یا سازمان، داده‌هایی با توزیع آماری خاص خود جمع‌آوری می‌کند. این توزیع‌ها می‌توانند با یکدیگر تفاوت چشمگیری داشته باشند، که در ادبیات این حوزه از آن با عنوان توزیع \gls{NonIID} یاد می‌شود.
    
    \tick \textbf{گره‌های معیوب یا مخرب}: در یک محیط ابری چندگانه واقعی، برخی گره‌ها ممکن است دچار خرابی نرم‌افزاری شوند، خط مشی‌های اشتباه داشته باشند یا در بدترین حالت، به‌طور عمدی به‌روزرسانی‌های مخرب ارسال کنند.
\end{itemize}

\gls{FedAvg} هیچ مکانیزمی برای مقابله با هیچ‌یک از این مشکلات ندارد و در حضور آن‌ها عملکردش به شدت افت می‌کند.

\section{چالش‌ها و انگیزه}
\label{sec:motivation}

روش‌های موجود برای مقاوم‌سازی یادگیری فدرال عمدتاً یکی از دو رویکرد کلی را دنبال می‌کنند. دسته اول روش‌هایی نظیر \lr{Krum} \cite{blanchard2017machine} هستند که به‌جای تجمیع همه به‌روزرسانی‌ها، تنها گره‌هایی را که به یکدیگر نزدیک‌ترند انتخاب می‌کنند. دسته دوم روش‌هایی نظیر میانه وزن‌دار \cite{yin2018byzantine} و میانگین کوتاه‌شده \cite{yin2018byzantine} هستند که به‌صورت مختصه‌به‌مختصه عمل می‌کنند.

هر دو دسته با چالش‌های جدی روبه‌رو هستند. \lr{Krum} پیچیدگی محاسباتی $O(n^2 d)$ دارد که در محیط‌های ابری با هزینه مستقیم محاسباتی همراه است. روش‌های مبتنی بر میانه نیز اطلاعات ارزشمند گره‌های صادق را از دست می‌دهند و سرعت همگرایی را کاهش می‌دهند. از سوی دیگر، هیچ‌یک از این روش‌ها ناهمگونی داده و مشکل گره‌های کُند را به‌صورت یکپارچه حل نمی‌کنند.

در محیط‌های ابری واقعی، یک مشکل اضافه نیز وجود دارد: یک گره ممکن است به دلیل یک اتفاق گذرا مثل مکث کوتاه ابری یا نوسان موقت شبکه، در یک دور ارتباطی خاص عملکرد بدی داشته باشد، اما در دورهای بعدی کاملاً طبیعی رفتار کند. اگر این گره به‌طور دائمی کنار گذاشته شود، اطلاعات ارزشمند داده‌های محلی آن از دست می‌رود.

این ملاحظات انگیزه اصلی این پژوهش را شکل می‌دهند: طراحی یک مکانیزم اجماع که بتواند به‌صورت تطبیقی و با حافظه تاریخی، رفتار هر گره را ارزیابی کند و وزن آن در تجمیع را بر اساس این ارزیابی تعیین کند؛ همه اینها بدون آنکه هزینه محاسباتی قابل‌توجهی به سیستم اضافه شود.


\section{نوآوری‌های پژوهش}
\label{sec:contributions}

نوآوری‌های اصلی این پایان‌نامه به شرح زیر است:

\begin{itemize}
    \tick 
    \textbf{ارائه مکانیزم \lr{TWAC}}: یک مکانیزم اجماع سبک‌وزن با پیچیدگی $O(nd)$ که از امتیازات اعتماد تاریخی برای وزن‌دهی به گره‌ها استفاده می‌کند. این روش بر پایه دو مؤلفه همسویی جهتی و منطقی‌بودن بزرگی طراحی شده تا در برابر حملات جهت‌معکوس‌سازی، مقیاس‌افزایی و نویز تصادفی مقاوم باشد؛ کارایی این روش به‌طور تجربی در برابر حمله جهت‌معکوس‌سازی ارزیابی شده است.

    \tick 
    \textbf{مدل‌سازی یکپارچه ناهمگونی}: بستری که ناهمگونی محاسباتی، شبکه‌ای و داده را به‌طور یکپارچه در یک پیاده‌سازی شبیه‌سازی می‌کند و برای اولین بار امکان مقایسه منصفانه روش‌های مختلف را در این فضای ناهمگون فراهم می‌آورد.

    \tick 
    \textbf{مدیریت گره‌های کُند}: مکانیزمی برای مدیریت گره‌هایی که به دلیل محدودیت‌های شبکه‌ای یا محاسباتی با تأخیر پاسخ می‌دهند، به‌گونه‌ای که تأخیر گذرا موجب محرومیت دائمی گره از فرایند یادگیری نشود.

    \tick 
    \textbf{ارزیابی تجربی جامع}: مقایسه کامل روش پیشنهادی با \lr{FedAvg}، میانگین کوتاه‌شده و \lr{Multi-Krum} در شش سناریو آزمایشی متفاوت که ترکیب‌های مختلفی از ناهمگونی داده، حمله بیزانتین و گره‌های کُند را پوشش می‌دهند.
\end{itemize}

\section{ساختار گزارش}
\label{sec:structure}

ساختار این پایان‌نامه به شرح زیر سازمان‌دهی شده است. نخست در 
\autoref{chap:concepts}،
تعاریف و مفاهیم پایه‌ای در حوزه یادگیری فدرال، پردازش توزیع‌شده و بسترهای ابری بیان می‌شود. در
\autoref{chap:relatedworks}،
کارهای پیشین در حوزه‌های مقاوم‌سازی یادگیری فدرال، مدیریت ناهمگونی و روش‌های اجماع مورد بررسی قرار می‌گیرند. در
\autoref{chap:approach}،
روش پیشنهادی این پژوهش یعنی \gls{TWAC} به‌تفصیل تشریح می‌شود که شامل مدل‌سازی ریاضی، الگوریتم پیشنهادی و تحلیل پیچیدگی است. در
\autoref{chap:simulation}،
بستر شبیه‌سازی، آزمایش‌های انجام‌شده و تحلیل نتایج ارائه می‌گردد و در نهایت در
\autoref{chap:conclusion}،
نتیجه‌گیری و پیشنهادات برای کارهای آینده بیان می‌شود.
